\clearpage\pagestyle{otMain}\fontsize{11}{15}\selectfont\corrigesChapitre{10}{Statistiques et boîtes à moustaches}
\begin{blocLivre}{cor-10-1}\begin{corrige}{1}{Trois indicateurs}\demarche{Additionner, repérer le rang central, puis soustraire les extrêmes.}
La somme vaut 55 : la moyenne vaut $55/5=\res{11}$. La valeur centrale, de rang 3, donne la médiane $\res{10}$. L’étendue vaut $15-7=\res8$. Moyenne et médiane répondent à des questions différentes.
\end{corrige}\end{blocLivre}
\begin{blocLivre}{cor-10-2}\begin{corrige}{2}{Une fréquence}\demarche{Identifier le total de référence.}
La fréquence est $12/30=\res{0{,}4}$, soit $\res{40\,\%}$. Le dénominateur est l’effectif de tout le groupe.
\end{corrige}\end{blocLivre}
\begin{blocLivre}{cor-10-3}\begin{corrige}{3}{Des effectifs cumulés}\demarche{Ajouter chaque effectif au cumul précédent.}
On additionne les effectifs de proche en proche : $\res{2\,;5\,;11\,;16\,;20}$. L’effectif total est le dernier effectif cumulé, $\res{20}$. Par exemple, 11 élèves ont une note inférieure ou égale à 10.
\end{corrige}\end{blocLivre}
\begin{blocLivre}{cor-10-4}\begin{corrige}{4}{Ordonner avant la médiane}\demarche{Ranger avant de choisir le rang central.}
La première liste ordonnée est $4\,;7\,;9\,;11\,;13$ : la médiane est la 3\textsuperscript{e} valeur, $\res9$. La seconde a six valeurs : on retient $(6+8)/2=\res7$. La valeur 20 n’influence pas cette médiane.
\end{corrige}\end{blocLivre}
\begin{blocLivre}{cor-10-5}\begin{corrige}{5}{Lire une boîte à moustaches}\demarche{Associer chaque trait de la boîte à un indicateur.}
On lit $\res{4\,;8\,;11\,;14\,;19}$ pour le minimum, $Q_1$, la médiane, $Q_3$ et le maximum. L’étendue vaut $19-4=\res{15\text{ min}}$. Comme $Q_3=14$, au moins $\res{75\,\%}$ des durées sont inférieures ou égales à 14 min.
\label{fig-cor-10-5}\figureGeo[corrige]{boite-lecture}
\end{corrige}\end{blocLivre}
\begin{blocLivre}{cor-10-6}\begin{corrige}{6}{Moyenne et quartiles d’un tableau}\demarche{Utiliser les effectifs cumulés pour trouver les rangs.}
La somme pondérée vaut $12+24+60+60+64=220$ : moyenne $\res{11}$. Avec $N=20$, les rangs 10 et 11 portent 10 : médiane $\res{10}$. $Q_1$ est la 5\textsuperscript{e} valeur, d’effectif cumulé atteint à 8 : $\res{Q_1=8}$. $Q_3$ est la 15\textsuperscript{e}, atteinte à 12 : $\res{Q_3=12}$.
\end{corrige}\end{blocLivre}
\begin{blocLivre}{cor-10-7}\begin{corrige}{7}{Quartiles d’une série brute}\demarche{Calculer les rangs, arrondis à l’entier supérieur.}
$N=11$. La médiane est la 6\textsuperscript{e} valeur : $\res{10}$. $11\times0{,}25=2{,}75$, donc $Q_1$ est la 3\textsuperscript{e} valeur : $\res6$. $11\times0{,}75=8{,}25$, donc $Q_3$ est la 9\textsuperscript{e} valeur : $\res{15}$.
\end{corrige}\end{blocLivre}
\begin{blocLivre}{cor-10-8}\begin{corrige}{8}{Construire une boîte}\demarche{Placer les cinq valeurs avant de dessiner la boîte.}
On place le minimum 3, $Q_1=6$, la médiane 10, $Q_3=15$ et le maximum 20. La boîte va de 6 à 15, avec un trait à 10 ; les moustaches vont de 3 à 6 et de 15 à 20. Le corrigé la montre en vert, avec l’échelle de l’axe.
\label{fig-cor-10-8}\figureGeo[corrige]{boite-construire}
\end{corrige}\end{blocLivre}
\begin{blocLivre}{cor-10-9}\begin{corrige}{9}{Quartiles et diagramme en barres}\demarche{Passer des hauteurs aux effectifs cumulés.}
Les effectifs 3, 7, 9, 6, 5 donnent les cumuls 3, 10, 19, 25, 30 : $\res{N=30}$. Les rangs 15 et 16 valent 3 livres : médiane $\res3$. $30\times0{,}25=7{,}5$ : rang 8, $\res{Q_1=2}$. $30\times0{,}75=22{,}5$ : rang 23, $\res{Q_3=4}$.
\label{fig-cor-10-9}\figureGeo[corrige]{stat-barres}
\end{corrige}\end{blocLivre}
\begin{blocLivre}{cor-10-10}\begin{corrige}{10}{Des formules de tableur}\demarche{Désigner une plage de cellules, puis contrôler à la main.}
On peut saisir \code{=MOYENNE(A2:A9)}, \code{=MEDIANE(A2:A9)} et \code{=MAX(A2:A9)-MIN(A2:A9)}. La moyenne vaut $114/8=\res{14{,}25}$. La série ordonnée est 9, 11, 12, 14, 15, 15, 18, 20 : médiane $(14+15)/2=\res{14{,}5}$. L’étendue vaut $\res{11}$.
\end{corrige}\end{blocLivre}
\begin{blocLivre}{cor-10-11}\begin{corrige}{11}{Réunir deux groupes}\demarche{Reconstituer les sommes de chaque groupe.}
Les sommes valent $12\times13=156$ et $18\times11=198$. La moyenne commune vaut $(156+198)/30=\res{11{,}8}$. Faire $(13+11)/2$ donnerait un poids égal à des groupes de tailles différentes.
\end{corrige}\end{blocLivre}
\begin{blocLivre}{cor-10-12}\begin{corrige}{12}{Comparer deux collèges}\demarche{Comparer d’abord les médianes, puis les longueurs des boîtes et les étendues.}
Les médianes sont 15 min pour A et 14 min pour B : positions proches. A est plus dispersé : étendue $40-5=35$ contre $25-8=17$, boîte de longueur $22-10=12$ contre $17-12=5$. Dans $\res{A}$, $Q_3=22$ : moins de 75 \% des durées sont inférieures à 22 min, donc au moins 25 \% des élèves mettent 22 min ou plus. Dans B, $Q_3=17$ : on sait seulement qu’au plus un quart des élèves mettent plus de 17 min, on ne peut pas l’affirmer.
\label{fig-cor-10-12}\figureGeo[corrige]{boites-comparer}
\end{corrige}\end{blocLivre}
\begin{blocLivre}{cor-10-13}\begin{corrige}{13}{Ajouter une observation}\demarche{Séparer ce que détermine la somme de ce qui demande l’ordre.}
L’ancienne somme est 126 ; la nouvelle vaut 150 pour 10 valeurs. La moyenne est $\res{15}$. La médiane reste inconnue : la somme ne renseigne pas sur les valeurs centrales. Par exemple neuf fois 14, ou huit fois 13 et une fois 22, ont la même moyenne initiale mais des médianes différentes après l’ajout.
\end{corrige}\end{blocLivre}
\begin{blocLivre}{cor-10-14}\begin{corrige}{14}{Même moyenne, mêmes résultats ?}\demarche{Comparer un indicateur de position et un de dispersion.}
Les deux sommes valent 40 : les moyennes sont 10. Les deux médianes sont 10. Les étendues sont $\res2$ et $\res{16}$. B est beaucoup plus dispersée ; les seuls indicateurs de position masquent cet écart.
\end{corrige}\end{blocLivre}
\begin{blocLivre}{cor-10-15}\begin{corrige}{15}{Un graphique trompeur}\demarche{Ramener le pourcentage à la valeur initiale.}
La hausse réelle est $(84-80)/80=\res{5\,\%}$. Les hauteurs visibles sont 2 et 6, dans un rapport 3. Ce rapport visuel ne traduit pas un triplement des inscriptions. Lire l’origine et les graduations avant de comparer.
\label{fig-cor-10-15}\figureGeo[corrige]{stat-axe-tronque}
\end{corrige}\end{blocLivre}
